\section{复数的概念}

\begin{review}
    对于二次方程：$x^2+x+1=0$ 的解应该如何表达？
\end{review}

\v{6em}

\begin{review}
    当 $a, b, c$ 都是实数且 $a \neq 0$ 时，关于 $x$ 的方程 $a x^2+b x+c=0$ 在复数范围内总是有解的，且：
    \begin{itemize}
        \item $\Delta=b^2-4 a c>0$ 时，方程有两个不相等的实数根；
        \item $\Delta=b^2-4 a c=0$ 时，方程有两个相等的实数根；
        \item $\Delta=b^2-4 a c<0$ 时，方程有两个互为共轭的虚数根．
    \end{itemize}
\end{review}

\begin{definition}[复数的概念]
    \[
        \text{复数}
        \begin{cases}
            \text{实数}
            \begin{cases}
                \text{有理数}
                \begin{cases}
                    \text{整数}
                    \begin{cases}
                        \text{正整数} \\
                        0          \\
                        \text{负整数}
                    \end{cases} \\
                    \text{分数}
                \end{cases} \\
                \text{无理数}
            \end{cases} \\
            \text{虚数}
            \begin{cases}
                \text{纯虚数} \\
                \text{非纯虚数}
            \end{cases}
        \end{cases}
    \]
    形如 $a+bi(a,b\in \mathbb{R})$ 的数叫做复数, 其中$i$叫做虚数单位, 满足 $i^{2}=-1$; $a$ 称为实部, $b$ 称为虚部（$bi$ 不是虚部）。
    \begin{tips}
        在上海，我们会使用 $\mathbf{Im}$ 表示虚部，使用 $\mathbf{Re}$ 表示实部。实际上是单词 Imaginary 和 Real 的缩写。
    \end{tips}
\end{definition}

\begin{definition}[复数的分类]
    \begin{minipage}{0.4\textwidth}
        对于复数 $z=a+bi(a,b\in \mathbb{R})$:
        \begin{itemize}
            \item $z$ 为实数 $\Leftrightarrow b=0$;
            \item $z$ 为虚数 $\leftrightarrow b \ne 0$;
            \item $z$ 为纯虚数 $\Leftrightarrow \begin{cases}a=0\\ b\ne0\end{cases}.$
        \end{itemize}
    \end{minipage}
    \begin{minipage}{0.58\textwidth}
        \begin{warning}
            若 $a=0,b=0$ 则 $z= 0$，即为实数.
        \end{warning}
    \end{minipage}
\end{definition}

\begin{example}
    设 $m \in \mathbf{R}, m^2+m-2+\left(m^2-1\right) \mathrm{i}$ 是纯虚数，则 $m=$ \underline{\qquad}.
\end{example}

\v{4em}
